% Standalone problem document, generated from the adjacent README.md.
% Compile with XeLaTeX. Mathematical content is maintained in Markdown.
\documentclass[11pt,a4paper]{article}
\usepackage[fontsize=10.5pt]{scrextend}
\usepackage[margin=22mm,headheight=15pt,headsep=9mm,footskip=11mm]{geometry}
\usepackage{fontspec}
\setmainfont{texgyrepagella}[Extension=.otf,UprightFont=*-regular,BoldFont=*-bold,ItalicFont=*-italic,BoldItalicFont=*-bolditalic]
\setsansfont{texgyreheros}[Extension=.otf,UprightFont=*-regular,BoldFont=*-bold,ItalicFont=*-italic,BoldItalicFont=*-bolditalic]
\setmonofont{texgyrecursor}[Extension=.otf,UprightFont=*-regular,BoldFont=*-bold,ItalicFont=*-italic,BoldItalicFont=*-bolditalic,Scale=MatchLowercase]
\usepackage{amsmath,amssymb}
\usepackage{unicode-math}
\setmathfont{texgyrepagella-math.otf}
\usepackage{xcolor,microtype,enumitem,needspace,fancyhdr}
\usepackage{bookmark}
\definecolor{ink}{HTML}{17344B}
\definecolor{accent}{HTML}{197D80}
\definecolor{muted}{HTML}{596773}
\hypersetup{colorlinks=true,linkcolor=accent,urlcolor=accent,pdftitle={IE-17 - Monotonic
optimal backward error along
LSMR},pdfauthor={Open Problems in Numerical Linear Algebra}}
\urlstyle{same}
\setlength{\parindent}{0pt}
\setlength{\parskip}{5pt plus 1pt minus 1pt}
\linespread{1.04}
\setlength{\emergencystretch}{3em}
\widowpenalty=10000
\clubpenalty=10000
\displaywidowpenalty=10000
\allowdisplaybreaks[1]
\setlist{nosep,leftmargin=1.5em}
\providecommand{\tightlist}{\setlength{\itemsep}{0pt}\setlength{\parskip}{0pt}}
\providecommand{\pandocbounded}[1]{#1}
\pagestyle{fancy}
\fancyhf{}
\fancyhead[L]{\sffamily\footnotesize\color{muted}OPEN PROBLEMS IN NUMERICAL LINEAR ALGEBRA}
\fancyhead[R]{\sffamily\bfseries\footnotesize\color{accent}IE-17}
\fancyfoot[L]{\parbox[b]{0.9\textwidth}{\sffamily\fontsize{8}{10}\selectfont\color{muted}Editorial ratings; a bounded search cannot certify that a problem remains unsolved.\\Literature check: 2026-09-08}}
\fancyfoot[R]{\sffamily\footnotesize\color{muted}\thepage}
\renewcommand{\headrulewidth}{0.3pt}
\renewcommand{\headrule}{\hbox to\headwidth{\color{accent}\leaders\hrule height \headrulewidth\hfill}}
\makeatletter
\renewcommand{\section}{\Needspace{5\baselineskip}\@startsection{section}{1}{0pt}{12pt plus 2pt minus 2pt}{4pt}{\sffamily\bfseries\large\color{ink}}}
\renewcommand{\subsection}{\Needspace{4\baselineskip}\@startsection{subsection}{2}{0pt}{9pt plus 1pt minus 1pt}{3pt}{\sffamily\bfseries\normalsize\color{ink}}}
\makeatother
\setcounter{secnumdepth}{0}
\begin{document}
{\sffamily\small\color{muted}Linear systems and elimination\par}
\vspace{3pt}
{\raggedright\hyphenpenalty=10000\exhyphenpenalty=10000\sffamily\bfseries\fontsize{21}{24}\selectfont\color{ink}Monotonic
optimal backward error along LSMR\par}
\vspace{7pt}
{\sffamily\small\textbf{Difficulty:} challenging\quad\textbf{Importance:} interesting
to the community\par}
\vspace{4pt}
{\color{accent}\hrule height 0.8pt}
\vspace{6pt}
\textbf{Status:} explicit dissertation conjecture; later resolution not
located

Let \(A\in\mathbb R^{m\times n}\) and \(b\in\mathbb R^m\). In exact
arithmetic, start LSMR at \(x_0=0\): equivalently, \(x_k\) minimizes
\(\|A^T(b-Ax)\|_2\) over \(\mathcal K_k(A^TA,A^Tb)\). Work up to its
exact termination and use its minimum-length iterate if necessary. Let
\(r=b-Ax\) and define the matrix-only normwise backward error \[
\mu(x)=\min\{\|E\|_2:(A+E)^T((A+E)x-b)=0\}.
\] For \(x\ne0\), put \(K_x=[A^T,\,(\|r\|_2/\|x\|_2)I]^T\),
\(v_x=[r^T,0^T]^T\), and \[
\widetilde\mu(x)=\frac{\|K_xK_x^{\dagger}v_x\|_2}{\|x\|_2}.
\] Here \(\dagger\) is the Moore--Penrose inverse; at an exact
least-squares solution set both errors to zero.

Are both sequences \(\mu(x_k)\) and \(\widetilde\mu(x_k)\)
nonincreasing, for successive nonzero LSMR iterates? These two closely
related claims are counted together, as in the source. The right-hand
side \(b\) is kept fixed in the backward-error model.

A positive answer would justify backward-error stopping decisions
without a later iteration making the current iterate less backward
accurate. The known monotonicity of \(\|r_k\|_2\) and \(\|A^Tr_k\|_2\)
does not establish this statement.

\subsection{References}\label{references}

D. C.-L. Fong, \emph{Minimum-Residual Methods for Sparse Least-Squares
Using Golub--Kahan Bidiagonalization}, Stanford dissertation 2011,
§7.2.1 and definitions (4.1),(4.5), printed pp.~56--57,118
(\href{https://web.stanford.edu/group/SOL/dissertations/david-fong-thesis-online.pdf}{primary
PDF}). Fong and M. Saunders, \emph{LSMR: An iterative algorithm for
sparse least-squares problems}, SISC 33(2011),2950--2971, §§6.1--6.2,
Fig7.5
(\href{https://web.stanford.edu/group/SOL/software/lsmr/LSMR-SISC-2011.pdf}{author
PDF}). E. Hallman and M. Gu, \emph{LSMB: Minimizing the backward error
for least-squares problems}, SIMAX 39(2018),1295--1317
(\href{https://erhallma.math.ncsu.edu/papers/hallman2018lsmb.pdf}{author
PDF}).

\subsection{Status check --- 2026-09-08}\label{status-check-2026-09-08}

Exact-title and ``LSMR optimal backward error monotonicity
conjecture/proof/counterexample'' searches found no resolution. The
LSMB2018 discussion of nonmonotonicity concerns LSQR, and its proposed
method differs from LSMR. Hallman's May 2026 arXiv:2605.09211 develops
backward-error formulas and estimates, without resolving this
iterate-monotonicity conjecture. No recent explicit reaffirmation was
located.
\end{document}
